\subsection{Summary of Contributions}
\label{sec:conclusion-contributions}

This paper makes four primary contributions to laser ultrasonic non-destructive testing:

\begin{enumerate}
\item \textbf{PMDF-Net multi-domain fusion architecture.} A dual-branch encoder-decoder that jointly processes time-domain waveforms and time-frequency representations, extracting complementary features for enhanced signal denoising fidelity.

\item \textbf{Physics-constrained loss function.} A feature-preserving loss formulation that enforces arrival-time, amplitude, and waveform-morphology consistency, ensuring denoised signals retain diagnostic information for defect characterization.

\item \textbf{Reduced averaging training protocol.} A paired training strategy using limited-average inputs ($N=16$) against high-average targets ($M=256$), enabling practical deployment without requiring extensive signal averaging during inspection.

\item \textbf{Cross-domain generalization.} Demonstrated transfer from intact aluminum specimens to defect-containing workpieces and across thickness variations, validating the approach for real-world inspection scenarios beyond training conditions.
\end{enumerate}

These contributions establish a framework for efficient laser ultrasonic inspection that reduces averaging requirements while preserving---and in some cases exceeding---conventional averaging performance in defect detection tasks.
